\documentclass[../../main.tex]{subfiles}

\begin{document}

\subsection{Il fascio di neutrini} \label{subsec:il-fascio-di-neutrini}
  \begin{figure}[ht]
    \includegraphics[width=\textwidth]{DUNE/dune-project/lbnf}
    \caption{Illustrazione della ``\textit{Long-Baseline Neutrino Facility}''
    del Fermilab e del complesso del ``\textit{Near Detector}''}
    \label{fig:lbnf}
  \end{figure}
  I detector di DUNE saranno esposti al fascio di neutrini più intenso del
  mondo, fornito dalla ``\textit{Long-Baseline Neutrino Facility}'' (LBNF)
  del Fermilab.\\
  Un fascio di protoni in un ampio spettro energetico di 60--120GeV viene
  estratto dal sincrotrone ``\textit{Main Injector}'' (MI) del Fermilab e
  trasportato verso l'area del bersaglio, costituito da un
  insieme lungo e sottile di segmenti di grafite in cui circa l'85\% dei
  protoni interagisce producendo mesoni carichi, prevalentemente
  pioni $\Ppipm$ e kaoni $\PKpm$.
  Le particelle secondarie entrano poi in un tubo di decadimento lungo
  194 m e riempito di elio, dove una grande frazione dei mesoni decadrà in
  $\Pmuon + \APnum$ o $\APmuon + \Pnum$ in base al segno del mesone, creando il
  fascio di neutrini.
  Il segno dei mesoni, e quindi la predominanza di $\Pnum$ o $\APnum$, è
  selezionato tramite due corna magnetiche (\textit{magnetic horns})
  posizionate prima del tubo di decadimento; le due modalità operative sono
  dette rispettivamente ``\textit{Forward Horn Current}'' (FHC) e
  ``\textit{Reverse Horn Current}'' (RHC).
  Il tubo di decadimento è seguito immediatamente dall'assorbitore, una
  struttura composta da alluminio, acciaio e cemento, con il compito di
  fermare i protoni che non hanno interagito nel bersaglio e i mesoni che non
  sono decaduti in neutrini.
  Una sezione longitudinale della linea di fascio LBNF è mostrata in figura
  ~\ref{fig:lbnf}.

  Questa struttura di focalizzazione è ottimizzata per fornire un fascio di
  neutrini a banda larga con un intervallo energetico tra 0,5 e 5 GeV, così
  da coprire il primo e il secondo massimo di oscillazione dei neutrini che,
  per una \textit{baseline} di 1300 km, si trovano approssimativamente a 2,4
  e 0,8 GeV\@.

  Il fascio ottenuto con questo processo è composto per circa il 90\% di
  $\Pnum$~($\APnum$), le impurità sono in parte costituite da antiparticelle del
  neutrino muonico desiderato e da alcuni $\APnue$~($\Pnue$), infatti i
  $\Pmuon$~($\APmuon$) possono decadere in $\Pelectron$~($\APelectron$) e
  $\APnue$~($\Pnue$), la maggior parte delle impurità restanti è costituita
  dai prodotti di decadimento dei kaoni poiché c'è la possibilità che questi
  decadano in neutrini (antineutrini) elettronici e muonici.

  La LBNF è progettata per operare a 1,2 MW o a 2,4 MW; l'esperimento si
  svolgerà nell'arco di 20 anni, durante i primi 5 la \textit{beam line}
  opererà a 1,2 MW, per poi salire a 2,4 MW per i restanti 15.\\
  Il fascio di neutrini sarà diviso in \textit{spill}, che alla potenza di 1,
  2 MW avranno una durata di \SI{e-5}{\second}, intervallati da pause di
  \SI{1,2}{\second} e saranno generati da \num{7,5e15} ``protoni su
  bersaglio'' (\textit{Proton On Target}) (POT).

\end{document}